\documentclass[12pt,a4paper]{article}
\usepackage[utf8]{vietnam}
\usepackage{amsmath,amssymb}
\usepackage{tabularx}
\usepackage{tikz}
\usepackage{xcolor}
\usepackage[top=0.6cm, bottom=0.6cm, left=0.6cm, right=0.6cm]{geometry}
\usepackage[dethi]{ex_test}

\SAOPTN{widthSA=3.2,heightSA=0.62}

\setlength{\parskip}{0pt}

% Dòng kẻ chấm mờ đều đặn chuẩn viết tay (khoảng cách 6.5mm rộng rãi vừa vặn)
\newcommand{\dongke}[1]{%
  \begingroup%
  \setlength{\parskip}{0pt}%
  \baselineskip=18.5pt%
  \vspace{1.0mm}%
  \foreach \i in {1,...,#1}{%
    \noindent{\color{gray!60}\dotfill}\par%
  }%
  \vspace{0.6mm}%
  \endgroup%
}

% Định nghĩa ô bài tập tự luận trong bảng
\newcommand{\oBaiTap}[3][6]{%
  \vspace{1.0mm}%
  \noindent\textbf{#2} #3\par%
  \dongke{#1}%
}

\begin{document}
\pagestyle{empty}
% =============================================================================
% ========================= TOÁN 6 - TUẦN 06 - BUỔI 01: MÃ ĐỀ 101 =============
% ========================= MẶT TRƯỚC (TRANG 1/2) ==============================
% =============================================================================

\noindent
\begin{tabularx}{\linewidth}{|X|c|}
\hline
\textbf{PHIẾU KIỂM TRA ĐỊNH KỲ TOÁN 6 - TUẦN 06 - BUỔI 01} & \textbf{Mã đề: 101 - Trang 1/2} \\
\hline
\end{tabularx}

\vspace{1mm}

\noindent
\begin{tabularx}{\linewidth}{|p{3.2cm}|X|}
\hline
\centering\textbf{ĐIỂM SỐ} & \centering\textbf{LỜI NHẬN XÉT CỦA THẦY CÔ} \tabularnewline
\hline
\rule{0pt}{1.8cm} & \tabularnewline
\hline
\end{tabularx}

\vspace{1.2mm}

\noindent\textbf{Câu 1 (0,25 điểm).} Một hình vuông có độ dài cạnh bằng $9\text{ cm}$. Chu vi của hình vuông đó bằng bao nhiêu? \shortans{36\text{ cm}}\par\vspace{0.3mm}
\noindent\textbf{Câu 2 (0,25 điểm).} Một mảnh đất hình chữ nhật có chiều dài $12\text{ m}$ và chiều rộng $5\text{ m}$. Diện tích của mảnh đất đó bằng bao nhiêu? \shortans{60\text{ m}^2}\par\vspace{0.3mm}
\noindent\textbf{Câu 3 (0,25 điểm).} Một hình thoi có độ dài hai đường chéo lần lượt là $8\text{ cm}$ và $10\text{ cm}$. Diện tích của hình thoi đó bằng bao nhiêu? \shortans{40\text{ cm}^2}\par\vspace{0.3mm}
\noindent\textbf{Câu 4 (0,25 điểm).} Một hình bình hành có độ dài cạnh đáy là $12\text{ cm}$ và chiều cao tương ứng là $6\text{ cm}$. Diện tích của hình bình hành đó bằng bao nhiêu? \shortans{72\text{ cm}^2}\par\vspace{0.3mm}
\noindent\textbf{Câu 5 (0,25 điểm).} Giá trị của biểu thức $2^4 \cdot 5 - 3^2 \cdot 4$ bằng bao nhiêu? \shortans{44}\par\vspace{0.3mm}
\noindent\textbf{Câu 6 (0,25 điểm).} Một hình vuông có diện tích bằng $64\text{ cm}^2$. Độ dài cạnh của hình vuông đó bằng bao nhiêu? \shortans{8\text{ cm}}\par\vspace{0.3mm}
\noindent\textbf{Câu 7 (0,25 điểm).} Một hình chữ nhật có chiều dài $15\text{ cm}$ và chiều rộng $8\text{ cm}$. Chu vi của hình chữ nhật đó bằng bao nhiêu? \shortans{46\text{ cm}}\par\vspace{0.3mm}
\noindent\textbf{Câu 8 (0,25 điểm).} Một hình thoi có chu vi bằng $36\text{ cm}$. Độ dài cạnh của hình thoi đó bằng bao nhiêu? \shortans{9\text{ cm}}\par\vspace{0.3mm}
\noindent\textbf{Câu 9 (0,25 điểm).} Một hình bình hành có diện tích bằng $48\text{ cm}^2$ và độ dài đáy bằng $8\text{ cm}$. Chiều cao tương ứng bằng bao nhiêu? \shortans{6\text{ cm}}\par\vspace{0.3mm}
\noindent\textbf{Câu 10 (0,25 điểm).} Giá trị của biểu thức $5 \cdot 2^3 - 3^3$ bằng bao nhiêu? \shortans{13}\par\vspace{0.3mm}
\vspace{0.4mm}

\noindent\textbf{Câu 11 (2,0 điểm). Thực hiện phép tính (hợp lý nếu có thể):}
\vspace{0.6mm}

\noindent
\begin{tabularx}{\linewidth}{|X|X|}
\hline
\oBaiTap[8]{a)}{$135 - 3^2 \cdot 10 + 2^3 \cdot 5$} & 
\oBaiTap[8]{b)}{$47 \cdot 68 + 47 \cdot 32 - 700$} \\
\hline
\oBaiTap[8]{c)}{$160 : [65 - (18 - 11)^2]$} & 
\oBaiTap[8]{d)}{$200 : \{ 80 - [22 + (9 - 7)^3] \}$} \\
\hline
\end{tabularx}

\newpage
% =============================================================================
% ========================= TOÁN 6 - TUẦN 06 - BUỔI 01: MÃ ĐỀ 101 =============
% ========================= MẶT SAU (TRANG 2/2) ================================
% =============================================================================

\noindent
\begin{tabularx}{\linewidth}{|X|c|}
\hline
\textbf{TOÁN 6 - TUẦN 06 - BUỔI 01} & \textbf{Mã đề: 101 - Trang 2/2} \\
\hline
\end{tabularx}

\vspace{1.5mm}

\noindent\textbf{Câu 12 (2,0 điểm). Tìm số tự nhiên $x$, biết:}
\vspace{0.8mm}

\noindent
\begin{tabularx}{\linewidth}{|X|X|}
\hline
\oBaiTap[5]{a)}{$45 + 5 \cdot (x - 3) = 75$} & 
\oBaiTap[5]{b)}{$95 - (3x + 14) = 45$} \\
\hline
\oBaiTap[5]{c)}{$2 \cdot 3^{x+1} - 15 = 39$} & 
\oBaiTap[5]{d)}{Một hình chữ nhật có chu vi bằng $40\text{ cm}$, chiều rộng $x\text{ cm}$, dài $12\text{ cm}$. Tìm $x$.} \\
\hline
\end{tabularx}

\vspace{1.5mm}

\noindent
\begin{tabularx}{\linewidth}{|X|}
\hline
\vspace{0.5mm}
\textbf{Câu 13 (2,5 điểm).} Bác An có một mảnh vườn hình chữ nhật có chiều dài $20\text{ m}$ và chiều rộng $12\text{ m}$.\par
\textbf{a) (1,25 điểm)} Tính chu vi và diện tích của mảnh vườn đó.\par
\textbf{b) (1,25 điểm)} Bác An dự định làm hàng rào lưới thép bao quanh khu vườn, chừa lại một cổng ra vào rộng $2\text{ m}$. Tính chiều dài hàng rào lưới thép và số tiền bác An phải trả để mua lưới thép, biết giá mỗi mét lưới thép là $60\ 000$ đồng.\par
\vspace{0.4mm}
\noindent\textit{Bài giải:}\dongke{11}
\\ \hline
\end{tabularx}

\vspace{1.5mm}

\noindent
\begin{tabularx}{\linewidth}{|X|}
\hline
\vspace{0.5mm}
\textbf{Câu 14 (1,0 điểm).} Cho tổng $A = 1 + 3 + 3^2 + 3^3 + \dots + 3^{99} + 3^{100}$. Chứng minh rằng $2A + 1$ là một lũy thừa của $3$.\par
\vspace{0.4mm}
\noindent\textit{Bài giải:}\dongke{7}
\\ \hline
\end{tabularx}

\newpage

% =============================================================================
% ========================= TOÁN 6 - TUẦN 06 - BUỔI 01: MÃ ĐỀ 102 =============
% ========================= MẶT TRƯỚC (TRANG 1/2) ==============================
% =============================================================================

\noindent
\begin{tabularx}{\linewidth}{|X|c|}
\hline
\textbf{PHIẾU KIỂM TRA ĐỊNH KỲ TOÁN 6 - TUẦN 06 - BUỔI 01} & \textbf{Mã đề: 102 - Trang 1/2} \\
\hline
\end{tabularx}

\vspace{1mm}

\noindent
\begin{tabularx}{\linewidth}{|p{3.2cm}|X|}
\hline
\centering\textbf{ĐIỂM SỐ} & \centering\textbf{LỜI NHẬN XÉT CỦA THẦY CÔ} \tabularnewline
\hline
\rule{0pt}{1.8cm} & \tabularnewline
\hline
\end{tabularx}

\vspace{1.2mm}

\noindent\textbf{Câu 1 (0,25 điểm).} Một hình vuông có độ dài cạnh bằng $8\text{ cm}$. Chu vi của hình vuông đó bằng bao nhiêu? \shortans{32\text{ cm}}\par\vspace{0.3mm}
\noindent\textbf{Câu 2 (0,25 điểm).} Một mảnh đất hình chữ nhật có chiều dài $15\text{ m}$ và chiều rộng $4\text{ m}$. Diện tích của mảnh đất đó bằng bao nhiêu? \shortans{60\text{ m}^2}\par\vspace{0.3mm}
\noindent\textbf{Câu 3 (0,25 điểm).} Một hình thoi có độ dài hai đường chéo lần lượt là $6\text{ cm}$ và $12\text{ cm}$. Diện tích của hình thoi đó bằng bao nhiêu? \shortans{36\text{ cm}^2}\par\vspace{0.3mm}
\noindent\textbf{Câu 4 (0,25 điểm).} Một hình bình hành có độ dài cạnh đáy là $15\text{ cm}$ và chiều cao tương ứng là $4\text{ cm}$. Diện tích của hình bình hành đó bằng bao nhiêu? \shortans{60\text{ cm}^2}\par\vspace{0.3mm}
\noindent\textbf{Câu 5 (0,25 điểm).} Giá trị của biểu thức $3^3 \cdot 2 - 2^4 \cdot 3$ bằng bao nhiêu? \shortans{6}\par\vspace{0.3mm}
\noindent\textbf{Câu 6 (0,25 điểm).} Một hình vuông có diện tích bằng $81\text{ cm}^2$. Độ dài cạnh của hình vuông đó bằng bao nhiêu? \shortans{9\text{ cm}}\par\vspace{0.3mm}
\noindent\textbf{Câu 7 (0,25 điểm).} Một hình chữ nhật có chiều dài $16\text{ cm}$ và chiều rộng $7\text{ cm}$. Chu vi của hình chữ nhật đó bằng bao nhiêu? \shortans{46\text{ cm}}\par\vspace{0.3mm}
\noindent\textbf{Câu 8 (0,25 điểm).} Một hình thoi có chu vi bằng $48\text{ cm}$. Độ dài cạnh của hình thoi đó bằng bao nhiêu? \shortans{12\text{ cm}}\par\vspace{0.3mm}
\noindent\textbf{Câu 9 (0,25 điểm).} Một hình bình hành có diện tích bằng $54\text{ cm}^2$ và độ dài đáy bằng $9\text{ cm}$. Chiều cao tương ứng bằng bao nhiêu? \shortans{6\text{ cm}}\par\vspace{0.3mm}
\noindent\textbf{Câu 10 (0,25 điểm).} Giá trị của biểu thức $4 \cdot 3^2 - 2^4$ bằng bao nhiêu? \shortans{20}\par\vspace{0.3mm}
\vspace{0.4mm}

\noindent\textbf{Câu 11 (2,0 điểm). Thực hiện phép tính (hợp lý nếu có thể):}
\vspace{0.6mm}

\noindent
\begin{tabularx}{\linewidth}{|X|X|}
\hline
\oBaiTap[8]{a)}{$145 - 2^3 \cdot 10 + 3^2 \cdot 5$} & 
\oBaiTap[8]{b)}{$53 \cdot 74 + 53 \cdot 26 - 800$} \\
\hline
\oBaiTap[8]{c)}{$180 : [54 - (16 - 10)^2]$} & 
\oBaiTap[8]{d)}{$300 : \{ 90 - [22 + (8 - 6)^3] \}$} \\
\hline
\end{tabularx}

\newpage
% =============================================================================
% ========================= TOÁN 6 - TUẦN 06 - BUỔI 01: MÃ ĐỀ 102 =============
% ========================= MẶT SAU (TRANG 2/2) ================================
% =============================================================================

\noindent
\begin{tabularx}{\linewidth}{|X|c|}
\hline
\textbf{TOÁN 6 - TUẦN 06 - BUỔI 01} & \textbf{Mã đề: 102 - Trang 2/2} \\
\hline
\end{tabularx}

\vspace{1.5mm}

\noindent\textbf{Câu 12 (2,0 điểm). Tìm số tự nhiên $x$, biết:}
\vspace{0.8mm}

\noindent
\begin{tabularx}{\linewidth}{|X|X|}
\hline
\oBaiTap[5]{a)}{$38 + 4 \cdot (x - 2) = 62$} & 
\oBaiTap[5]{b)}{$86 - (4x + 10) = 44$} \\
\hline
\oBaiTap[5]{c)}{$3 \cdot 2^{x+1} - 12 = 36$} & 
\oBaiTap[5]{d)}{Một hình thoi có chu vi bằng $52\text{ cm}$ và độ dài cạnh bằng $x\text{ cm}$. Tìm $x$.} \\
\hline
\end{tabularx}

\vspace{1.5mm}

\noindent
\begin{tabularx}{\linewidth}{|X|}
\hline
\vspace{0.5mm}
\textbf{Câu 13 (2,5 điểm).} Bác Bình có một mảnh vườn hình chữ nhật có chiều dài $25\text{ m}$ và chiều rộng $10\text{ m}$.\par
\textbf{a) (1,25 điểm)} Tính chu vi và diện tích của mảnh vườn đó.\par
\textbf{b) (1,25 điểm)} Bác Bình dự định làm hàng rào lưới thép bao quanh khu vườn, chừa lại một cổng ra vào rộng $3\text{ m}$. Tính chiều dài hàng rào lưới thép và số tiền bác Bình phải trả để mua lưới thép, biết giá mỗi mét lưới thép là $55\ 000$ đồng.\par
\vspace{0.4mm}
\noindent\textit{Bài giải:}\dongke{11}
\\ \hline
\end{tabularx}

\vspace{1.5mm}

\noindent
\begin{tabularx}{\linewidth}{|X|}
\hline
\vspace{0.5mm}
\textbf{Câu 14 (1,0 điểm).} Cho tổng $B = 1 + 2 + 2^2 + 2^3 + \dots + 2^{99} + 2^{100}$. Chứng minh rằng $B + 1$ là một lũy thừa của $2$.\par
\vspace{0.4mm}
\noindent\textit{Bài giải:}\dongke{7}
\\ \hline
\end{tabularx}

\newpage

% =============================================================================
% ========================= TOÁN 6 - TUẦN 06 - BUỔI 01: MÃ ĐỀ 103 =============
% ========================= MẶT TRƯỚC (TRANG 1/2) ==============================
% =============================================================================

\noindent
\begin{tabularx}{\linewidth}{|X|c|}
\hline
\textbf{PHIẾU KIỂM TRA ĐỊNH KỲ TOÁN 6 - TUẦN 06 - BUỔI 01} & \textbf{Mã đề: 103 - Trang 1/2} \\
\hline
\end{tabularx}

\vspace{1mm}

\noindent
\begin{tabularx}{\linewidth}{|p{3.2cm}|X|}
\hline
\centering\textbf{ĐIỂM SỐ} & \centering\textbf{LỜI NHẬN XÉT CỦA THẦY CÔ} \tabularnewline
\hline
\rule{0pt}{1.8cm} & \tabularnewline
\hline
\end{tabularx}

\vspace{1.2mm}

\noindent\textbf{Câu 1 (0,25 điểm).} Một hình vuông có độ dài cạnh bằng $7\text{ cm}$. Chu vi của hình vuông đó bằng bao nhiêu? \shortans{28\text{ cm}}\par\vspace{0.3mm}
\noindent\textbf{Câu 2 (0,25 điểm).} Một mảnh đất hình chữ nhật có chiều dài $14\text{ m}$ và chiều rộng $5\text{ m}$. Diện tích của mảnh đất đó bằng bao nhiêu? \shortans{70\text{ m}^2}\par\vspace{0.3mm}
\noindent\textbf{Câu 3 (0,25 điểm).} Một hình thoi có độ dài hai đường chéo lần lượt là $8\text{ cm}$ và $14\text{ cm}$. Diện tích của hình thoi đó bằng bao nhiêu? \shortans{56\text{ cm}^2}\par\vspace{0.3mm}
\noindent\textbf{Câu 4 (0,25 điểm).} Một hình bình hành có độ dài cạnh đáy là $14\text{ cm}$ và chiều cao tương ứng là $5\text{ cm}$. Diện tích của hình bình hành đó bằng bao nhiêu? \shortans{70\text{ cm}^2}\par\vspace{0.3mm}
\noindent\textbf{Câu 5 (0,25 điểm).} Giá trị của biểu thức $2^5 \cdot 3 - 5^2 \cdot 2$ bằng bao nhiêu? \shortans{46}\par\vspace{0.3mm}
\noindent\textbf{Câu 6 (0,25 điểm).} Một hình vuông có diện tích bằng $49\text{ cm}^2$. Độ dài cạnh của hình vuông đó bằng bao nhiêu? \shortans{7\text{ cm}}\par\vspace{0.3mm}
\noindent\textbf{Câu 7 (0,25 điểm).} Một hình chữ nhật có chiều dài $18\text{ cm}$ và chiều rộng $6\text{ cm}$. Chu vi của hình chữ nhật đó bằng bao nhiêu? \shortans{48\text{ cm}}\par\vspace{0.3mm}
\noindent\textbf{Câu 8 (0,25 điểm).} Một hình thoi có chu vi bằng $44\text{ cm}$. Độ dài cạnh của hình thoi đó bằng bao nhiêu? \shortans{11\text{ cm}}\par\vspace{0.3mm}
\noindent\textbf{Câu 9 (0,25 điểm).} Một hình bình hành có diện tích bằng $63\text{ cm}^2$ và độ dài đáy bằng $7\text{ cm}$. Chiều cao tương ứng bằng bao nhiêu? \shortans{9\text{ cm}}\par\vspace{0.3mm}
\noindent\textbf{Câu 10 (0,25 điểm).} Giá trị của biểu thức $6 \cdot 2^3 - 4^2$ bằng bao nhiêu? \shortans{32}\par\vspace{0.3mm}
\vspace{0.4mm}

\noindent\textbf{Câu 11 (2,0 điểm). Thực hiện phép tính (hợp lý nếu có thể):}
\vspace{0.6mm}

\noindent
\begin{tabularx}{\linewidth}{|X|X|}
\hline
\oBaiTap[8]{a)}{$150 - 4^2 \cdot 5 + 2^4 \cdot 3$} & 
\oBaiTap[8]{b)}{$36 \cdot 83 + 36 \cdot 17 - 600$} \\
\hline
\oBaiTap[8]{c)}{$210 : [51 - (15 - 9)^2]$} & 
\oBaiTap[8]{d)}{$240 : \{ 85 - [37 + (7 - 5)^3] \}$} \\
\hline
\end{tabularx}

\newpage
% =============================================================================
% ========================= TOÁN 6 - TUẦN 06 - BUỔI 01: MÃ ĐỀ 103 =============
% ========================= MẶT SAU (TRANG 2/2) ================================
% =============================================================================

\noindent
\begin{tabularx}{\linewidth}{|X|c|}
\hline
\textbf{TOÁN 6 - TUẦN 06 - BUỔI 01} & \textbf{Mã đề: 103 - Trang 2/2} \\
\hline
\end{tabularx}

\vspace{1.5mm}

\noindent\textbf{Câu 12 (2,0 điểm). Tìm số tự nhiên $x$, biết:}
\vspace{0.8mm}

\noindent
\begin{tabularx}{\linewidth}{|X|X|}
\hline
\oBaiTap[5]{a)}{$52 + 6 \cdot (x - 4) = 88$} & 
\oBaiTap[5]{b)}{$98 - (5x + 13) = 45$} \\
\hline
\oBaiTap[5]{c)}{$4 \cdot 3^{x-1} - 10 = 26$} & 
\oBaiTap[5]{d)}{Một hình bình hành có diện tích bằng $84\text{ cm}^2$, chiều cao $7\text{ cm}$, cạnh đáy tương ứng là $x\text{ cm}$. Tìm $x$.} \\
\hline
\end{tabularx}

\vspace{1.5mm}

\noindent
\begin{tabularx}{\linewidth}{|X|}
\hline
\vspace{0.5mm}
\textbf{Câu 13 (2,5 điểm).} Bác Cường có một mảnh vườn hình chữ nhật có chiều dài $18\text{ m}$ và chiều rộng $12\text{ m}$.\par
\textbf{a) (1,25 điểm)} Tính chu vi và diện tích của mảnh vườn đó.\par
\textbf{b) (1,25 điểm)} Bác Cường dự định làm hàng rào lưới thép bao quanh khu vườn, chừa lại một cổng ra vào rộng $2\text{ m}$. Tính chiều dài hàng rào lưới thép và số tiền bác Cường phải trả để mua lưới thép, biết giá mỗi mét lưới thép là $65\ 000$ đồng.\par
\vspace{0.4mm}
\noindent\textit{Bài giải:}\dongke{11}
\\ \hline
\end{tabularx}

\vspace{1.5mm}

\noindent
\begin{tabularx}{\linewidth}{|X|}
\hline
\vspace{0.5mm}
\textbf{Câu 14 (1,0 điểm).} Cho tổng $C = 1 + 4 + 4^2 + 4^3 + \dots + 4^{49} + 4^{50}$. Chứng minh rằng $3C + 1$ là một lũy thừa của $4$.\par
\vspace{0.4mm}
\noindent\textit{Bài giải:}\dongke{7}
\\ \hline
\end{tabularx}

\newpage

% =============================================================================
% ========================= TOÁN 6 - TUẦN 06 - BUỔI 01: MÃ ĐỀ 104 =============
% ========================= MẶT TRƯỚC (TRANG 1/2) ==============================
% =============================================================================

\noindent
\begin{tabularx}{\linewidth}{|X|c|}
\hline
\textbf{PHIẾU KIỂM TRA ĐỊNH KỲ TOÁN 6 - TUẦN 06 - BUỔI 01} & \textbf{Mã đề: 104 - Trang 1/2} \\
\hline
\end{tabularx}

\vspace{1mm}

\noindent
\begin{tabularx}{\linewidth}{|p{3.2cm}|X|}
\hline
\centering\textbf{ĐIỂM SỐ} & \centering\textbf{LỜI NHẬN XÉT CỦA THẦY CÔ} \tabularnewline
\hline
\rule{0pt}{1.8cm} & \tabularnewline
\hline
\end{tabularx}

\vspace{1.2mm}

\noindent\textbf{Câu 1 (0,25 điểm).} Một hình vuông có độ dài cạnh bằng $11\text{ cm}$. Chu vi của hình vuông đó bằng bao nhiêu? \shortans{44\text{ cm}}\par\vspace{0.3mm}
\noindent\textbf{Câu 2 (0,25 điểm).} Một mảnh đất hình chữ nhật có chiều dài $16\text{ m}$ và chiều rộng $5\text{ m}$. Diện tích của mảnh đất đó bằng bao nhiêu? \shortans{80\text{ m}^2}\par\vspace{0.3mm}
\noindent\textbf{Câu 3 (0,25 điểm).} Một hình thoi có độ dài hai đường chéo lần lượt là $10\text{ cm}$ và $12\text{ cm}$. Diện tích của hình thoi đó bằng bao nhiêu? \shortans{60\text{ cm}^2}\par\vspace{0.3mm}
\noindent\textbf{Câu 4 (0,25 điểm).} Một hình bình hành có độ dài cạnh đáy là $18\text{ cm}$ và chiều cao tương ứng là $5\text{ cm}$. Diện tích của hình bình hành đó bằng bao nhiêu? \shortans{90\text{ cm}^2}\par\vspace{0.3mm}
\noindent\textbf{Câu 5 (0,25 điểm).} Giá trị của biểu thức $3^2 \cdot 10 - 2^3 \cdot 5$ bằng bao nhiêu? \shortans{50}\par\vspace{0.3mm}
\noindent\textbf{Câu 6 (0,25 điểm).} Một hình vuông có diện tích bằng $100\text{ cm}^2$. Độ dài cạnh của hình vuông đó bằng bao nhiêu? \shortans{10\text{ cm}}\par\vspace{0.3mm}
\noindent\textbf{Câu 7 (0,25 điểm).} Một hình chữ nhật có chiều dài $20\text{ cm}$ và chiều rộng $6\text{ cm}$. Chu vi của hình chữ nhật đó bằng bao nhiêu? \shortans{52\text{ cm}}\par\vspace{0.3mm}
\noindent\textbf{Câu 8 (0,25 điểm).} Một hình thoi có chu vi bằng $56\text{ cm}$. Độ dài cạnh của hình thoi đó bằng bao nhiêu? \shortans{14\text{ cm}}\par\vspace{0.3mm}
\noindent\textbf{Câu 9 (0,25 điểm).} Một hình bình hành có diện tích bằng $72\text{ cm}^2$ và độ dài đáy bằng $8\text{ cm}$. Chiều cao tương ứng bằng bao nhiêu? \shortans{9\text{ cm}}\par\vspace{0.3mm}
\noindent\textbf{Câu 10 (0,25 điểm).} Giá trị của biểu thức $7 \cdot 2^3 - 5^2$ bằng bao nhiêu? \shortans{31}\par\vspace{0.3mm}
\vspace{0.4mm}

\noindent\textbf{Câu 11 (2,0 điểm). Thực hiện phép tính (hợp lý nếu có thể):}
\vspace{0.6mm}

\noindent
\begin{tabularx}{\linewidth}{|X|X|}
\hline
\oBaiTap[8]{a)}{$160 - 5^2 \cdot 4 + 3^3 \cdot 2$} & 
\oBaiTap[8]{b)}{$64 \cdot 69 + 64 \cdot 31 - 900$} \\
\hline
\oBaiTap[8]{c)}{$240 : [52 - (18 - 12)^2]$} & 
\oBaiTap[8]{d)}{$280 : \{ 95 - [47 + (6 - 4)^3] \}$} \\
\hline
\end{tabularx}

\newpage
% =============================================================================
% ========================= TOÁN 6 - TUẦN 06 - BUỔI 01: MÃ ĐỀ 104 =============
% ========================= MẶT SAU (TRANG 2/2) ================================
% =============================================================================

\noindent
\begin{tabularx}{\linewidth}{|X|c|}
\hline
\textbf{TOÁN 6 - TUẦN 06 - BUỔI 01} & \textbf{Mã đề: 104 - Trang 2/2} \\
\hline
\end{tabularx}

\vspace{1.5mm}

\noindent\textbf{Câu 12 (2,0 điểm). Tìm số tự nhiên $x$, biết:}
\vspace{0.8mm}

\noindent
\begin{tabularx}{\linewidth}{|X|X|}
\hline
\oBaiTap[5]{a)}{$41 + 7 \cdot (x - 3) = 76$} & 
\oBaiTap[5]{b)}{$105 - (6x + 15) = 42$} \\
\hline
\oBaiTap[5]{c)}{$5 \cdot 2^{x-1} - 14 = 26$} & 
\oBaiTap[5]{d)}{Một hình chữ nhật có diện tích bằng $96\text{ cm}^2$, chiều rộng bằng $8\text{ cm}$, chiều dài bằng $x\text{ cm}$. Tìm $x$.} \\
\hline
\end{tabularx}

\vspace{1.5mm}

\noindent
\begin{tabularx}{\linewidth}{|X|}
\hline
\vspace{0.5mm}
\textbf{Câu 13 (2,5 điểm).} Bác Dũng có một mảnh vườn hình chữ nhật có chiều dài $22\text{ m}$ và chiều rộng $10\text{ m}$.\par
\textbf{a) (1,25 điểm)} Tính chu vi và diện tích của mảnh vườn đó.\par
\textbf{b) (1,25 điểm)} Bác Dũng dự định làm hàng rào lưới thép bao quanh khu vườn, chừa lại một cổng ra vào rộng $2\text{ m}$. Tính chiều dài hàng rào lưới thép và số tiền bác Dũng phải trả để mua lưới thép, biết giá mỗi mét lưới thép là $50\ 000$ đồng.\par
\vspace{0.4mm}
\noindent\textit{Bài giải:}\dongke{11}
\\ \hline
\end{tabularx}

\vspace{1.5mm}

\noindent
\begin{tabularx}{\linewidth}{|X|}
\hline
\vspace{0.5mm}
\textbf{Câu 14 (1,0 điểm).} Cho tổng $D = 1 + 5 + 5^2 + 5^3 + \dots + 5^{49} + 5^{50}$. Chứng minh rằng $4D + 1$ là một lũy thừa của $5$.\par
\vspace{0.4mm}
\noindent\textit{Bài giải:}\dongke{7}
\\ \hline
\end{tabularx}

\newpage

\end{document}
